% Generated draft: requires mathematical and source review.
\documentclass[11pt]{article}
\usepackage[margin=25mm]{geometry}
\usepackage{fontspec}
\setmainfont{Latin Modern Roman}
\usepackage{amsmath,amssymb,mathtools,unicode-math}
\setmathfont{Latin Modern Math}
\usepackage{enumitem,needspace,xurl,hyperref}
\setlength{\parindent}{0pt}
\setlength{\parskip}{5pt}
\setlength{\emergencystretch}{3em}
\newcommand{\N}{\mathbb{N}}
\newcommand{\Z}{\mathbb{Z}}
\newcommand{\Q}{\mathbb{Q}}
\newcommand{\R}{\mathbb{R}}
\newcommand{\C}{\mathbb{C}}
\newcommand{\F}{\mathbb{F}}
\newcommand{\A}{\mathbb{A}}
\newcommand{\PP}{\mathbb{P}}
\newcommand{\E}{\mathbb{E}}
\newcommand{\eps}{\varepsilon}
\newcommand{\dd}{\,\mathrm{d}}
\begin{document}
\section*{UnsolvedMath ID: 10400214}
\textbf{Source problem code:} catalogue entry
\subsection{Definitions and mathematical statement}
A KTG is a knotted trivalent graph, with the framing or ribbon data and graph operations intended by the source. Problem 12.8 asks for a natural invariant of KTGs constructed from configuration-space integrals. For example, the ordinary ordered configuration space is
\[
 \operatorname{Conf}_k(X)=\{(x_1,\ldots,x_k)\in X^k:x_i\ne x_j\ (i\ne j)\};
\]
the requested construction must specify the relevant graph-constrained configurations, differential forms, integrations and any boundary corrections, and be invariant under the intended equivalences. The excerpt supplies none of those choices, a target space, or an exact meaning of naturality, and does not demand a particular named invariant.

The supplied subsequent material includes a separate task, Exercise 12.9: find a complete set of moves between shadow surfaces presenting homeomorphic \(3\)-manifolds. A shadow surface is the source's fake surface with an integer or half-integer gleam on each face. The excerpt describes a fake surface locally as an open subset of the cone over a tetrahedron; which skeleton of the tetrahedron is intended requires checking. It describes gleams as Chern numbers determining an appropriately defined \(S^1\)-bundle presentation and says every closed \(3\)-manifold admits such a closed-shadow presentation. Because the base can be singular, that appropriately defined bundle is not specified by the usual locally trivial bundle definition alone. The stated pentagon and hexagon moves preserve the presented manifold but are not complete; Exercise 12.9 asks for moves that connect any two presentations of the same manifold, with the intended homeomorphism conventions.

The context identifies a shadow as the time evolution of a sequence of KTG operations. Its Figure 22 relates framed links, their exteriors, surgery to closed \(3\)-manifolds, certain KTG sequences, collapsible shadow surfaces and capping their boundaries to closed shadows; compatibility of these presentations is attributed to [40, 384]. The shadow number of a manifold is the minimum number of vertices among its shadow presentations. The source attributes this notion to D. Thurston, motivated by [279, 272, 273], and reports shadow number \(0\) for graph manifolds and \(1\) for all Borromean-ring surgeries. The latter exact-value assertion needs qualification or source checking, rather than automatic replacement by an upper bound. The final reference to a following volume-related conjecture does not supply that conjecture's statement.
\subsection{Short English statement}
Construct a natural configuration-integral invariant of knotted trivalent graphs. The supplied continuation separately asks for complete shadow-presentation moves and supplies shadow-number background whose conventions require clarification.
\subsection{Sources}
\begin{itemize}
\item T. Ohtsuki, \emph{Problems on invariants of knots and 3-manifolds} (2002), Problem 12.8, following Exercise 12.9 and Figure 22, supplied PDF page 163. \url{https://msp.org/gtm/2002/04/gtm-2002-04-024s.pdf}. Recorded access 2026-07-29; extraction pdf-font-structure.
\item The excerpt attributes shadow presentations to Turaev [388], also referring to [40, 384], and cites [388, Figure 1.1 of Chapter VIII] for pentagon and hexagon relations. It attributes shadow number to D. Thurston with motivation from [279, 272, 273]. These numbered bibliography entries are not supplied.
\item Local UnsolvedMath registry, \texttt{data/unsolvedmath/problems.json}, permanent ID \(10400214\), exact source code \texttt{catalogue entry}; research classification OPEN-TRIAGE.
\end{itemize}
\end{document}
